\section{Conclusion}
\label{sec:conclusion}

We introduced a subject-calibrated EOG-guided diffusion method for ocular-artifact removal. The method estimates an EOG-to-EEG propagation matrix from a separate calibration segment, shrinks uncertain individual estimates toward population calibration, and uses the resulting ocular contribution to condition a shared diffusion denoiser. Matched calibration improved paired waveform reconstruction in every development participant and in seven of eight held-out participants, with similar mean effect sizes in the two cohorts. Natural recordings showed a complementary gain in EOG attenuation and low-EOG signal retention. A controlled injected sensitivity analysis was consistent with participant matching when calibration information was delivered through a support-fitted representation, and propagation-aware sample widths improved the empirical coverage of predictive intervals on development data.

The present scope is EOG-assisted, same-session ocular-artifact removal with a compatible electrode montage. Natural recordings provide attenuation and preservation measures rather than a clean waveform target, and the predictive-interval analysis remains a development result. Linear regression and deterministic networks were competitive point estimators, so the current evidence supports the complete calibration-guided system rather than a general point-error advantage of diffusion. The next decisive evaluations are an external EEG--EOG cohort, a compute-matched comparison between diffusion sampling and deterministic ensembles, and held-out assessment of the predictive intervals. Within the evaluated setting, the results show that participant adaptation can be achieved by calibrating the ocular propagation pathway while keeping the EEG model shared.
